% EEE3096S Practical 3, Voltage Mirror. Group 79.
% Every measured number lives in measurements.tex, never in this file.
% Captures go in figures/ as full12.png and coarse6.png.
%   make        build the PDF
%   make check  build, then list what is still unmeasured or missing
\documentclass[11pt,a4paper]{article}
\usepackage{lmodern}
\usepackage[margin=20mm]{geometry}
\usepackage{amsmath}
\usepackage{booktabs}
\usepackage{tabularx}
\usepackage{graphicx}
\usepackage{xcolor}
\usepackage[hidelinks]{hyperref}
\setlength{\parindent}{0pt}
\setlength{\parskip}{5pt}
\renewcommand{\arraystretch}{1.4}

% An unfilled measurement. Impossible to miss in a proof read.
\newcommand{\bench}[1]{\textcolor{red}{\textbf{[BENCH: #1]}}}

\newcommand{\captureplaceholder}[2]{%
  \fbox{\parbox[c][0.26\textheight][c]{\dimexpr\linewidth-2\fboxsep-2\fboxrule\relax}{%
    \centering \textcolor{red}{\textbf{#2}}\par\medskip
    Not captured yet. Save it as \texttt{\detokenize{#1}}.}}}
\newcommand{\capture}[2]{%
  \IfFileExists{#1}{%
    \includegraphics[width=\linewidth,height=0.26\textheight,keepaspectratio]{#1}}{%
    \captureplaceholder{#1}{#2}}}

% Every number in the report comes from here. A slot still holding \bench{}
% has not been measured yet, and `make check` lists it.
% Replace the whole \bench{...} with the reading, units out of the macro.

% Header
\newcommand{\mDate}{\bench{date}}
\newcommand{\mBuild}{\bench{Debug or Release}}

% Reference and predictions
\newcommand{\mVr}{\bench{VR}}
\newcommand{\mDmmRange}{\bench{DMM range and resolution}}
\newcommand{\mLsb}{\bench{VR/4096}}
\newcommand{\mCoarseStep}{\bench{VR/64}}

% DC row 1, target 0.800 V
\newcommand{\mLoFullIn}{\bench{PA5}}
\newcommand{\mLoFullC}{\bench{C}}
\newcommand{\mLoFullD}{\bench{D}}
\newcommand{\mLoFullPred}{\bench{D VR/4096}}
\newcommand{\mLoFullOut}{\bench{PA4}}
\newcommand{\mLoCoarseIn}{\bench{PA5}}
\newcommand{\mLoCoarseC}{\bench{C}}
\newcommand{\mLoCoarseD}{\bench{D}}
\newcommand{\mLoCoarsePred}{\bench{D VR/4096}}
\newcommand{\mLoCoarseOut}{\bench{PA4}}

% DC row 2, target 1.650 V
\newcommand{\mMidFullIn}{\bench{PA5}}
\newcommand{\mMidFullC}{\bench{C}}
\newcommand{\mMidFullD}{\bench{D}}
\newcommand{\mMidFullPred}{\bench{D VR/4096}}
\newcommand{\mMidFullOut}{\bench{PA4}}
\newcommand{\mMidCoarseIn}{\bench{PA5}}
\newcommand{\mMidCoarseC}{\bench{C}}
\newcommand{\mMidCoarseD}{\bench{D}}
\newcommand{\mMidCoarsePred}{\bench{D VR/4096}}
\newcommand{\mMidCoarseOut}{\bench{PA4}}

% DC row 3, target 2.500 V
\newcommand{\mHiFullIn}{\bench{PA5}}
\newcommand{\mHiFullC}{\bench{C}}
\newcommand{\mHiFullD}{\bench{D}}
\newcommand{\mHiFullPred}{\bench{D VR/4096}}
\newcommand{\mHiFullOut}{\bench{PA4}}
\newcommand{\mHiCoarseIn}{\bench{PA5}}
\newcommand{\mHiCoarseC}{\bench{C}}
\newcommand{\mHiCoarseD}{\bench{D}}
\newcommand{\mHiCoarsePred}{\bench{D VR/4096}}
\newcommand{\mHiCoarseOut}{\bench{PA4}}

% Triangle and step
\newcommand{\mGenFreq}{\bench{Hz}}
\newcommand{\mGenVpp}{\bench{Vpp}}
\newcommand{\mGenOffset}{\bench{offset}}
\newcommand{\mPaMin}{\bench{PA5 min}}
\newcommand{\mPaMax}{\bench{PA5 max}}
\newcommand{\mStepMeasured}{\bench{measured step}}
\newcommand{\mStepDelta}{\bench{measured minus predicted}}
\newcommand{\mStepPct}{\bench{percent}}


\begin{document}
\begin{center}
  {\Large\bfseries EEE3096S Practical 3 --- Voltage Mirror}\par
  Tri Tac Le, LXXTRI004 \quad Gil Sharon, SHRGIL004 \quad Group 79
\end{center}
Date: \mDate \qquad Build configuration: \mBuild

ADC channel 5 on PA5 is read at 2000 samples per second and the code is written
to DAC channel 1 on PA4. SW0 selects the mapping. Both converters are configured
for 12 bits and stay that way; COARSE6 discards the six least significant bits in
software only.

\section*{1. Measurements and predictions}

Measured analogue reference $V_R$: \mVr\par
DMM range and resolution: \mDmmRange\par
Nominal converter interval $V_R/4096$: \mLsb\par
Predicted coarse output step $V_R/64$: \mCoarseStep

PA5 and PA4 were measured independently with the DMM. $C$ and $D$ are the codes
shown on the LCD. The predicted output is $D\,V_R/4096$ using the recorded command
and the measured reference.

{\small
\begin{tabularx}{\linewidth}{@{}l l X r r X X@{}}
\toprule
Input target & Mode & PA5 (V) & $C$ & $D$ & Predicted PA4 (V) & PA4 (V) \\
\midrule
$\sim$0.800 V & FULL12  & \mLoFullIn   & \mLoFullC   & \mLoFullD   & \mLoFullPred   & \mLoFullOut \\
$\sim$0.800 V & COARSE6 & \mLoCoarseIn & \mLoCoarseC & \mLoCoarseD & \mLoCoarsePred & \mLoCoarseOut \\
$\sim$1.650 V & FULL12  & \mMidFullIn   & \mMidFullC   & \mMidFullD   & \mMidFullPred   & \mMidFullOut \\
$\sim$1.650 V & COARSE6 & \mMidCoarseIn & \mMidCoarseC & \mMidCoarseD & \mMidCoarsePred & \mMidCoarseOut \\
$\sim$2.500 V & FULL12  & \mHiFullIn   & \mHiFullC   & \mHiFullD   & \mHiFullPred   & \mHiFullOut \\
$\sim$2.500 V & COARSE6 & \mHiCoarseIn & \mHiCoarseC & \mHiCoarseD & \mHiCoarsePred & \mHiCoarseOut \\
\bottomrule
\end{tabularx}}

Generator: \mGenFreq, \mGenVpp, \mGenOffset\par
Verified PA5 minimum and maximum: \mPaMin\ and \mPaMax\par
Measured adjacent coarse step: \mStepMeasured\par
Measured minus predicted: \mStepDelta, which is \mStepPct

\newpage
\section*{2. Scope captures}

\begin{center}
\capture{figures/full12.png}{Capture 1, FULL12}
\end{center}
\textbf{Capture 1, FULL12.} PA5 on CH1 and PA4 on CH2, DC coupled, over about one
triangle cycle. The individual 12-bit steps are smaller than the scope resolves.

\begin{center}
\capture{figures/coarse6.png}{Capture 2, COARSE6}
\end{center}
\textbf{Capture 2, COARSE6.} The same two channels, zoomed until adjacent output
plateaus separate, with voltage cursors across two settled plateaus. The cursor
reading is the measured step above.

\newpage
\section*{3. Explanation}

\textbf{Mapping.} FULL12 is $D = C$. COARSE6 is $D = 64\lfloor C/64 \rfloor$,
implemented by clearing the six least significant bits of the code rather than by
changing either converter's resolution. That leaves 64 distinct commands, spaced
64 codes or \mCoarseStep\ apart. The largest is 4032 because 4032 is $63 \times 64$
and the next multiple of 64 is 4096, which is past the highest 12-bit code of 4095.
There are 4096 codes but the highest code is 4095, which is why the count and the
maximum differ.

\textbf{Measured steps.} The measured plateau spacing is \mStepMeasured\ against the
predicted $V_R/64 = \mCoarseStep$, a difference of \mStepDelta. The prediction uses
the nominal model $V = D\,V_R/4096$, so it carries the reference measurement's own
uncertainty, and the cursor can only be placed to the scope's vertical resolution on
a trace that still has noise on it. Buffer offset and the loading of the probe move
the absolute level more than they move the spacing between two neighbouring
plateaus, which is why the step compares better than any single absolute reading.

\textbf{Resolution and accuracy.} The FULL12 trace looks continuous because its step
is $V_R/4096 = \mLsb$, below what the scope resolves on this timebase and smaller
than the noise on the trace, not because the output is continuous. It is still 4096
discrete levels. Resolution is only the size of the step; accuracy is whether the
step lands at the right voltage. A 12-bit converter with an offset, a gain error or
an uncertain reference still has 4096 levels and still produces the wrong voltage at
each one. Both converters share this board's reference, so a code copied from ADC to
DAC gives a nominally unity transfer and the shared reference cancels out of the
ratio. That cancellation says nothing about either converter's absolute accuracy.

\textbf{One DC result.} At the 0.800 V target the ADC estimate is
$C\,V_R/4096$ from $C = \mLoFullC$ against a measured input of \mLoFullIn, and the
nominal output is $D\,V_R/4096 = \mLoFullPred$ from $D = \mLoFullD$ against a
measured output of \mLoFullOut. Copying the code does not guarantee the same
physical voltage: the ADC quantises the input to the nearest code and loses
everything inside that interval, and the DAC then adds its own offset, gain error
and output buffer behaviour on the way back out. The code is the same number at both
ends, but the two conversions are separate physical operations with separate errors.

\end{document}
